\documentclass[10pt,a4paper]{amsart}
\usepackage[margin=19mm]{geometry}
\usepackage{amsmath,amssymb,mathtools}
\usepackage[scr=rsfs,cal=euler]{mathalfa}
\usepackage{xfrac}
\usepackage[unicode,hidelinks,bookmarks=false]{hyperref}
\newcommand{\ie}{i.e.\ }
\newcommand{\cf}{cf.\ }
\newcommand{\resp}{respectively\ }
\newcommand{\Subsection}{\subsection}
\providecommand{\nobreakdash}{\nobreak\hbox{-}}
\setlength{\emergencystretch}{2em}
\multlinegap0pt
\usepackage{bm}
\usepackage{bbm}
\usepackage{dsfont}
\usepackage{xspace}
\usepackage{cancel}
\usepackage{mathtools}
\usepackage{amsthm}
\makeatletter
\@ifundefined{theorem}{\newtheorem{theorem}{Theorem}}{}
\@ifundefined{lemma}{\newtheorem{lemma}[theorem]{Lemma}}{}
\@ifundefined{proposition}{\newtheorem{proposition}[theorem]{Proposition}}{}
\makeatother
\newcommand{\N}{\mathbb N}% by default, $\N$ is American naturals $\{1,2,\ldots\}$
\renewcommand{\text}{\textup}
\title[On a Conjecture of Cusick on a sum of Cantor sets]{On a Conjecture of Cusick on a sum of Cantor sets}
\author{Nikita Shulga}
\date{Source: arXiv:2411.17379v2 (CC BY 4.0)}
\begin{document}
\maketitle

\noindent\emph{Source and licence.} On a Conjecture of Cusick on a sum of Cantor sets. Version arXiv:2411.17379v2 (CC BY 4.0), distributed under the Creative Commons Attribution 4.0 International licence, as recorded in the arXiv licence metadata for this exact version and shown on the abstract page https://arxiv.org/abs/2411.17379v2. Full rights notice: licenses/arxiv-2411.17379-CC-BY-4.0.txt.\par\smallskip

\noindent\textit{Editorial scope.} Counted unit: the Lemma whose source label is \texttt{bk\_correctly\_defined} in the article (source0.tex lines 881--883, source label \texttt{bk\_correctly\_defined}), delivered together with the article's own complete original proof of it (source0.tex lines 885--1036, one \texttt{proof} environment of 152 lines). No mathematics is written, repaired or re-derived: statement and proof are the article's own text line for line. Required context retained uncounted: range (lines 636--649); algorithm:cn (lines 771--773); condition:cylinder:x-pq (lines 828--830). Every definition, display and label of those lines is retained unchanged. Omitted from this package and not redistributed: 1--635, 650--770, 774--827, 831--880, 884--884, 1037--1900. Nothing inside a retained excerpt is omitted or altered: every retained body is delivered exactly as the article's lines are written, so no omission receipt of any kind accompanies this package. Citations: the retained excerpts carry no citation command at all, so no bibliography entry of the article is transcribed and no reference is redistributed. Reference adaptation: none was needed and none was made; every reference inside the delivered unit resolves inside it, so no unresolved reference or citation is produced. Printed slips kept literal: no printed word, symbol, hypothesis or number of the counted statement or of its proof was altered. Layout and class: the article's own class is \texttt{amsart} with packages including geometry, graphics, amsmath, amssymb, enumitem, comment, url, mathtools, graphicx, epsfig, hyperref, color, xcolor, pgf; the wrapper is the lane's pinned \texttt{amsart} editorial wrapper at 10pt on A4, which restates the theorem-like environments the retained text needs and adds the article's own macro definitions that the retained text needs (\texttt{\textbackslash N}, \texttt{\textbackslash text}), with extra wrapper packages (bm, bbm, dsfont, xspace, cancel, mathtools). The delivered numbering is the unit's own. Assets: no retained span contains a figure environment, a graphics-inclusion command, a PDF-page inclusion, a table or a TikZ picture; no image file is copied into this package, the package declares no asset dependency and no image file is redistributed with it. Licence: version arXiv:2411.17379v2 of this article is distributed under the Creative Commons Attribution 4.0 International licence, as recorded in the arXiv licence metadata for this exact version and shown on the abstract page https://arxiv.org/abs/2411.17379v2. A full rights notice is delivered at licenses/arxiv-2411.17379-CC-BY-4.0.txt. Characters: every retained excerpt is pure ASCII, so the delivered engine, XeLaTeX with its default Latin Modern text font, reports no missing character.
\begin{proposition}\label{range}
For any $n\ge 1$ and $\left(c_{1},\ldots,c_{n}\right)\in\mathbb{N}^{n}$, we have
\begin{equation}I_{n}\left(c_{1},\ldots,c_{n}\right)=\left\{
  \begin{array}{ll}
   \left[\dfrac{p_{n}}{q_{n}},\dfrac{p_{n}+p_{n-1}}{q_{n}+q_{n-1}}\right) &\textrm{if }n\textrm{ is even},\\
   \left(\dfrac{p_{n}+p_{n-1}}{q_{n}+q_{n-1}},\dfrac{p_{n}}{q_{n}}\right] &\textrm{if }n\textrm{ is odd}.\\
  \end{array}
\right.
\end{equation}
Therefore, the length of it is given by
\begin{equation}\label{length}
\left|I_{n}\left(c_{1},\ldots,c_{n}\right)\right|=\dfrac{1}{q_{n}\left(q_{n}+q_{n-1}\right)}.
\end{equation}
\end{proposition}

\begin{equation}\label{algorithm:cn}
c_{n+1} = a_{n+1} \left(    x - \frac{s_n}{t_n}  \right)  +1.
\end{equation}

\begin{equation}\label{condition:cylinder:x-pq}
a_i \left( x - \frac{p_n}{q_n} \right) = b_i  \text{ for } 1\leq i \leq n-1
\end{equation}

\begin{lemma}\label{bk_correctly_defined}
    Assume that $c_1,b_1,\ldots,b_{k-1},c_k$ are correctly defined. Then $b_k$ is also correctly defined.
\end{lemma}

\begin{proof}
Case 1. $k=2n$ is an even number. 

Inequality which defines $b_{2n}$ is 
\begin{equation}\label{deff:b2n}
\frac{s_{2n}}{t_{2n}} \leq x - \frac{p_{2n}}{q_{2n}} < \frac{s_{2n}+s_{2n-1}}{t_{2n}+t_{2n-1}}.
\end{equation}

The condition \eqref{condition:cylinder:x-pq} is equivalent to the fact that 
$$
x-\frac{p_{2n}}{q_{2n}} \in I_{2n-1}(b_1,\ldots,b_{2n-1}),
$$
which in turn by Proposition \ref{range} is equivalent to the inequality
\begin{equation}\label{induction:b2n}
\frac{s_{2n-1}+s_{2n-2}}{t_{2n-1}+t_{2n-2}} < x - \frac{p_{2n}}{q_{2n}} \leq \frac{s_{2n-1}}{t_{2n-1}}.
\end{equation}
We want to show that a number $x-p_{2n}/q_{2n}$ will always satisfy the inequality \eqref{induction:b2n}.

By inductive hypothesis, we know that
$$
\frac{p_{2n}-p_{2n-1}}{q_{2n}-q_{2n-1}} \leq x - \frac{s_{2n-1}}{t_{2n-1}} < \frac{p_{2n}}{q_{2n}},
$$
which can be rewritten as 
$$
\frac{s_{2n-1}}{t_{2n-1}} - \frac{1}{q_{2n}(q_{2n}-q_{2n-1})} \leq x - \frac{p_{2n}}{q_{2n}} < \frac{s_{2n-1}}{t_{2n-1}}.
$$
Thus right-hand side of inequality \eqref{induction:b2n} is automatically true. Now we need to deal with the left-hand side.

It is enough to show that
$$
\frac{s_{2n-1}}{t_{2n-1}} - \frac{1}{q_{2n}(q_{2n}-q_{2n-1})} > \frac{s_{2n-1}+s_{2n-2}}{t_{2n-1}+t_{2n-2}}
$$
or, rewritten using properties of neighbor Farey fractions,
\begin{equation}
\frac{1}{t_{2n-1}(t_{2n-1}+t_{2n-2})} > \frac{1}{q_{2n}(q_{2n}-q_{2n-1})},
\end{equation}
which is just
\begin{equation}\label{induction:b2n:final}
q_{2n}(q_{2n}-q_{2n-1}) >t_{2n-1}(t_{2n-1}+t_{2n-2}).
\end{equation}


Case 2.  $k=2n+1$ is an odd number. 

Inequality which defines $b_{2n+1}$ is 
\begin{equation}\label{deff:b2n+1}
\frac{s_{2n+1}+s_{2n}}{t_{2n+1}+t_{2n}} < x - \frac{p_{2n+1}}{q_{2n+1}} \leq    \frac{s_{2n+1}}{t_{2n+1}}.
\end{equation}

The condition \eqref{condition:cylinder:x-pq} is equivalent to the fact that 
$$
x-\frac{p_{2n+1}}{q_{2n+1}} \in I_{2n}(b_1,\ldots,b_{2n}),
$$
which in turn by Proposition \ref{range} is equivalent to the inequality
\begin{equation}\label{induction:b2n+1}
\frac{s_{2n}}{t_{2n}}   \leq x - \frac{p_{2n+1}}{q_{2n+1}} < \frac{s_{2n}+s_{2n-1}}{t_{2n}+t_{2n-1}}.
\end{equation}
We want to show that a number $x-p_{2n+1}/q_{2n+1}$ will always satisfy the inequality \eqref{induction:b2n+1}.

By inductive hypothesis, we know that
$$
\frac{p_{2n+1}}{q_{2n+1}}  < x - \frac{s_{2n}}{t_{2n}} \leq \frac{p_{2n+1}-p_{2n}}{q_{2n+1}-q_{2n}},
$$
which can be rewritten as 
$$
\frac{s_{2n}}{t_{2n}}  < x - \frac{p_{2n+1}}{q_{2n+1}} \leq \frac{s_{2n}}{t_{2n}}+ \frac{1}{q_{2n+1}(q_{2n+1}-q_{2n})}.
$$
Thus left-hand side of inequality \eqref{induction:b2n+1} is automatically true. Now we need to deal with the right-hand side.

It is enough to show that
$$
\frac{s_{2n}}{t_{2n}}+ \frac{1}{q_{2n+1}(q_{2n+1}-q_{2n})} < \frac{s_{2n}+s_{2n-1}}{t_{2n}+t_{2n-1}},
$$
or, rewritten using properties of neighbor Farey fractions,
\begin{equation}
\frac{1}{t_{2n}(t_{2n}+t_{2n-1})} > \frac{1}{q_{2n+1}(q_{2n+1}-q_{2n})},
\end{equation}
which is just
\begin{equation}\label{induction:b2n+1:final}
q_{2n+1}(q_{2n+1}-q_{2n}) >t_{2n}(t_{2n}+t_{2n-1}).
\end{equation}


As one can see, both \eqref{induction:b2n:final} and \eqref{induction:b2n+1:final} can be written in terms of $k$ as
\begin{equation}\label{induction:bnbn:finalefinale}
q_k(q_k-q_{k-1}) > t_{k-1}(t_{k-1}+t_{k-2}).
\end{equation}


To show this we need the following lemma.
 
\begin{lemma}\label{lemma:quotient:qlarge}
Assume that partial quotients $c_1,b_1,\ldots,b_{k-1},c_k$ are correctly defined. Then for all $1\leq n \leq k$ we have 
% $$
% \frac{q_{n+1}}{t_{n}} > \frac{t_{n}+t_{n-1}}{q_{n}} .
% $$
\begin{equation}\label{inequality:lemma1}
q_{n}q_{n-1} > t_{n-1}(t_{n-1}+t_{n-2}) .
\end{equation}
\end{lemma}

\begin{proof}
We will show this for the odd index $2m+1\leq k$. For even indices, the process is exactly the same except for inequality signs due to the properties of even and odd convergents. By assumption, all partial quotients $c_i$ with indices smaller or equal to $2m+1$ and all partial quotients $b_i$ with indices smaller or equal to $2k$ are correctly defined.

Thus by our procedure, $c_{2m+1}$ satisfies
$$
\frac{p_{2m+1}}{q_{2m+1}}  < x - \frac{s_{2m}}{t_{2m}}   \leq \frac{p_{2m+1}-p_{2m}}{q_{2m+1}-q_{2m}}.
$$
This implies
\begin{equation}\label{eq:good:first}
\frac{1}{q_{2m}q_{2m+1}}  < \left| x - \frac{s_{2m}}{t_{2m}} - \frac{p_{2m}}{q_{2m}} \right| \leq \frac{1}{q_{2m}(q_{2m+1}-q_{2m})}.
\end{equation}
On the other hand, by procedure we also have
$$
\frac{s_{2m}}{t_{2m}} \leq x- \frac{p_{2m}}{q_{2m}} < \frac{s_{2m}+s_{2m-1}}{t_{2m}+t_{2m-1}}.
$$
Thus
\begin{equation}\label{eq:good:second}
0 \leq \left| x - \frac{s_{2m}}{t_{2m}} - \frac{p_{2m}}{q_{2m}} \right| < \frac{1}{t_{2m}(t_{2m}+t_{2m-1})}.
\end{equation}
Combining left-hand side of \eqref{eq:good:first} with right-hand side of \eqref{eq:good:second} we get
$$
q_{2m}q_{2m+1} > t_{2m}(t_{2m}+t_{2m-1}).
$$
% or, in terms of $n$ this inequality is
% $$
% q_n q_{n-1} > t_{n-1}(t_{n-1}+t_{n-2})
% $$
% or, dividing both sides by $q_{2m}t_{2m}$,
% $$
% \frac{q_{2m+1}}{t_{2m}} > \frac{t_{2m}+t_{2m-1}}{q_{2m}}.
% $$
\end{proof}

\noindent
\textit{End of proof of Lemma \ref{bk_correctly_defined}. }
By Lemma \ref{lemma:quotient:qlarge} for $n=k$, we get
$$
q_{k}q_{k-1} > t_{k-1} (t_{k-1}+t_{k-2})
$$
Thus in order to secure \eqref{induction:bnbn:finalefinale} we need to check that
$$
q_{k}q_{k-1} < q_k(q_k-q_{k-1})
$$
The latter is equivalent to 
$$
2 q_{k-1} < q_{k}=c_{k} q_{k-1} + q_{k-2},
$$
which is always true as $c_i\geq2$ for all $i\in\N$ just by definition of procedure \eqref{algorithm:cn}.

Thus $b_k$ can always be correctly defined assuming that all previous partial quotients are correctly defined. 
\end{proof}

\end{document}
